\documentclass[10pt,a4paper]{amsart}
\usepackage[margin=19mm]{geometry}
\usepackage{amsmath,amssymb,mathtools}
\usepackage[scr=rsfs,cal=euler]{mathalfa}
\usepackage{xfrac}
\usepackage[unicode,hidelinks,bookmarks=false]{hyperref}
\newcommand{\ie}{i.e.\ }
\newcommand{\cf}{cf.\ }
\newcommand{\resp}{respectively\ }
\newcommand{\Subsection}{\subsection}
\providecommand{\nobreakdash}{\nobreak\hbox{-}}
\setlength{\emergencystretch}{2em}
\multlinegap0pt
\usepackage{amssymb}
\usepackage{enumerate}
\usepackage{xcolor}
\usepackage{tikz}
\newtheorem{lemma}{Lemma}
\title{A Weak Structural Form of Commutative Equivalence in Finite Codes}
\author{Dean Kraizberg}
\date{Source: arXiv:2603.27656v1 (CC BY 4.0)}
\begin{document}
\maketitle

\noindent\emph{Source and licence.} A Weak Structural Form of Commutative Equivalence in Finite Codes. Version arXiv:2603.27656v1 (CC BY 4.0), distributed by arXiv under the Creative Commons Attribution 4.0 International licence, http://creativecommons.org/licenses/by/4.0/, as recorded in the arXiv licence metadata for this exact version and shown on the abstract page https://arxiv.org/abs/2603.27656v1. Full licence notice: \texttt{licenses/arxiv-2603.27656-CC-BY-4.0.txt}.\par\smallskip

\noindent\textit{Editorial scope.} Counted unit: the article's lemma lemma:subset-sum-exact (source0.tex lines 552--558). The article's complete original proof of it is delivered with it (source0.tex lines 560--573, one \texttt{proof} environment of 14 lines). No mathematics is written, repaired or re-derived: the statement and the proof are the article's own lines, in the article's own notation, and no retained line is substituted or re-flowed.  No further source-local context is needed: every cross-reference of every retained line resolves inside the two delivered spans.  Nothing was omitted from any retained span, so the package carries no omission record.  No retained line contains a citation, so the wrapper carries no bibliography and no bibliography entry.  No retained line contains a figure environment, an image-inclusion command or drawing code, so no image is delivered and the package declares no asset dependency.  Editorial formatting only, in addition: an AMS \texttt{amsart} preamble that restates the article's own declarations and macros used by the retained lines, namely the environments \texttt{lemma} on separate counters, together with the standard packages those lines need.  The retained environments print in this document as Lemma 1 in order of appearance, whereas the article prints the same items with the numbering of its own full document.  The complete native source of the article as distributed in the arXiv e-print for this exact version is bundled unchanged as \texttt{sources/197370/source0.tex}.
\begin{lemma}\label{lemma:subset-sum-exact}
Let $A$ be a finite multiset of natural numbers, and let $N \in \mathbb{N}$. Suppose that 
\[
\sum_{n \in A} 2^n \ge 2^N \quad \text{and} \quad n \le N \text{ for all } n \in A.
\] 
Then there exists a sub-multiset $A' \subseteq A$ such that $\sum_{n \in A'} 2^n = 2^N$.
\end{lemma}

\begin{proof}
We proceed by induction on $N$. 

\textit{Base case:} $N = 0$.  
In this case, the only possible elements in $A$ are $0$, and the condition $\sum_{n \in A} 2^n \ge 2^0 = 1$ implies that $A$ contains at least one copy of $0$. Let $A'$ be a multiset consisting of exactly one $0$ from $A$. Then $\sum_{n \in A'} 2^n = 2^0 = 1$, as required.

\textit{Inductive step:} Consider $N > 0$, and assume that the statement holds for all natural numbers less than $N$. 

If $N \in A$, then we may take $A' = \{N\}$ and the result follows immediately. Otherwise, $N \notin A$. By assumption, $\sum_{n \in A} 2^n \ge 2^N \ge 2^{N-1}$. By the induction hypothesis applied to $N-1$, there exists a sub-multiset $A_1 \subseteq A$ such that $\sum_{n \in A_1} 2^n = 2^{N-1}.$ Consider the remaining multiset $A \setminus A_1$. We have 
\[
\sum_{n \in A \setminus A_1} 2^n = \sum_{n \in A} 2^n - \sum_{n \in A_1} 2^n \ge 2^N - 2^{N-1} = 2^{N-1}.
\] 
Applying the induction hypothesis again to $A \setminus A_1$ for $N-1$, we obtain a sub-multiset $A_2 \subseteq A \setminus A_1$ such that $\sum_{n \in A_2} 2^n = 2^{N-1}.$ Setting $A' = A_1 \cup A_2$, we have $\sum_{n \in A'} 2^n = 2^N$, as desired.
\end{proof}

\end{document}
